% Part: first-order-logic
% Chapter: models-theories
% Section: expressing-props-of-structures

\documentclass[../../../include/open-logic-section]{subfiles}

\begin{document}

\olfileid{fol}{mat}{exs}

\olsection{بیان خواص \printtoken{P}{structure}}

\begin{explain}
بیان شرایطی دربارهٔ توابع و روابط، یا به‌طور کلی‌تر بیان اینکه توابع و
روابط موجود در یک ساختار این شرایط را برآورده می‌کنند، اغلب سودمند و
مهم است. برای نمونه، می‌خواهیم راه‌هایی برای تمایز آن !!{structure}s
یک زبان که معنای مورد نظر ما برای !!{predicate}s را «بازمی‌نمایانند»
از ساختارهایی داشته باشیم که چنین نمی‌کنند. البته کاملاً آزادیم مشخص
کنیم کدام !!{structure}s را «مقصود» می‌دانیم؛ مثلاً می‌توانیم مقرر کنیم
که تعبیر !!{predicate}~$\le$ باید یک ترتیب باشد، یا تنها به تعبیرهایی
از~$\Lang L$ علاقه‌مند باشیم که در آن‌ها دامنه از مجموعه‌ها تشکیل شده و
$\Obj \in$ با رابطهٔ «!!a{element} از ... است» تعبیر می‌شود. اما آیا می‌توانیم این
کار را با !!{sentence}s زبان انجام دهیم؟ به بیان دیگر، کدام شرایط دربارهٔ
!!a{structure}~$\Struct M$ را می‌توانیم با !!a{sentence} (یا شاید
مجموعه‌ای از !!{sentence}s) در زبان~$\Struct M$ بیان کنیم؟ برخی شرایط را
نخواهیم توانست بیان کنیم. برای نمونه، هیچ جمله‌ای از~$\Lang L_A$
وجود ندارد که تنها در !!a{structure} چون~$\Struct M$ صادق باشد اگر
$\Domain M = \Nat$. نمی‌توانیم بیان کنیم «دامنه فقط شامل اعداد طبیعی
است». اما شاید بتوانیم «خواص ساختاری» !!{structure}s را بیان کنیم. کدام
خواص !!{structure}s را می‌توان با !!{sentence}s بیان کرد؟ یا به تعبیری
دیگر، کدام مجموعه‌های !!{structure}s را می‌توانیم ساختارهایی توصیف کنیم
که !!a{sentence} (یا مجموعه‌ای از !!{sentence}s) را صادق می‌سازند؟
\end{explain}

\begin{defn}[مدل یک مجموعه]
فرض کنید $\Gamma$ مجموعه‌ای از !!{sentence}s در زبان~$\Lang L$ باشد.
می‌گوییم !!a{structure}~$\Struct M$ \emph{مدلی از}~$\Gamma$ است اگر برای
هر $!A \in \Gamma$ داشته باشیم $\Sat{M}{!A}$.
\end{defn}

\begin{ex}
جملهٔ~$\lforall[x][x \le x]$ در~$\Struct M$ صادق است هرگاه و تنها
هرگاه $\Assign{\le}{M}$ رابطه‌ای بازتابی باشد. جملهٔ
$\lforall[x][\lforall[y][((x \le y \land y \le x) \lif x = y)]]$ در
~$\Struct M$ صادق است هرگاه و تنها هرگاه $\Assign{\le}{M}$ پادتقارن
باشد. جملهٔ~$\lforall[x][\lforall[y][\lforall[z][((x \le y \land y \le z)
      \lif x \le z)]]]$ در~$\Struct M$ صادق است هرگاه و تنها هرگاه
$\Assign{\le}{M}$ تراگذری باشد. بنابراین مدل‌های
\begin{align*}
\{\quad &\lforall[x][x \le x], \\
   & \lforall[x][\lforall[y][((x \le y \land y \le
    x) \lif x = y)]], \\
   &\lforall[x][\lforall[y][\lforall[z][((x \le y
      \land y \le z) \lif x \le z)]]] \quad \}
\end{align*}
دقیقاً ساختارهایی هستند که در آن‌ها~$\Assign{\le}{M}$ بازتابی،
پادتقارن و تراگذری، یعنی یک ترتیب جزئی، است. ازاین‌رو می‌توانیم این
این جمله‌ها را اصول موضوع \emph{نظریهٔ مرتبهٔ اول ترتیب‌های جزئی}
در نظر بگیریم.
\end{ex}

\end{document}
